% CHTC Researcher Forum 2026: Posters are 4 feet (48 inches) wide and 3 feet (36 inches) tall.
% Hack: Currently set in betterposter.cls
\documentclass[a0paper,fleqn]{betterposter}
\usepackage{bm}
\graphicspath{{figures/}}
\usepackage[hidelinks]{hyperref}

% Uncomment the following commands to customise the format

%% Setting the width of columns
% Left column
\setlength{\leftbarwidth}{0.2\paperwidth}
% Right column
\setlength{\rightbarwidth}{0.2\paperwidth}

%% Setting the column margins
% Horizontal margin
%\setlength{\columnmarginvertical}{0.05\paperheight}
% Vertical margin
%\setlength{\columnmarginhorizontal}{0.05\paperheight}
% Horizontal margin for the main column
%\setlength{\maincolumnmarginvertical}{0.15\paperheight}
% Vertical margin for the main column
%\setlength{\maincolumnmarginhorizontal}{0.15\paperheight}

% Text font
\renewcommand{\fontsizestandard}{\fontsize{32}{46} \selectfont}
% Main column font
% \renewcommand{\fontsizemain}{\fontsize{148.00}{280.00} \selectfont}
% Set baseline skip to 1.4 times font size
% \renewcommand{\fontsizemain}{\fontsize{135}{162} \selectfont}
% \renewcommand{\fontsizemain}{\fontsize{130}{156} \selectfont}
\renewcommand{\fontsizemain}{\fontsize{128}{180} \selectfont}
% \renewcommand{\fontsizemain}{\fontsize{130}{182} \selectfont}
% Title font
% \renewcommand{\fontsizetitle}{\fontsize{101.50}{152.00} \selectfont}
\renewcommand{\fontsizetitle}{\fontsize{64}{90} \selectfont}
% Author font
\renewcommand{\fontsizeauthor}{\fontsize{33.55}{47.3} \selectfont}
\newcommand{\fontsizeinstitution}{\fontsize{20}{25} \selectfont}
% Section font
\renewcommand{\fontsizesection}{\fontsize{42}{52} \selectfont}
\renewcommand{\fontsizesubsection}{\fontsize{36}{45} \selectfont}

% Changing font sizes for a specific text segment
% Place the text inside brackets:
% {\fontsize{28}{35} \selectfont Your text goes here}

%% Changing colours
% Background of side columns
%\renewcommand{\columnbackgroundcolor}{black}
% Font of side columns
%\renewcommand{\columnfontcolor}{gray}
% Background of main column
%\renewcommand{\maincolumnbackgroundcolor}{empirical}
%\renewcommand{\maincolumnbackgroundcolor}{theory}
%\renewcommand{\maincolumnbackgroundcolor}{methods}
%\renewcommand{\maincolumnbackgroundcolor}{intervention}
% Font of main column
%\renewcommand{\maincolumnfontcolor}{gray}

% Disable hyphenation
\hyphenpenalty=10000
\exhyphenpenalty=10000

\begin{document}
\betterposter{
 % MAIN COLUMN
 \maincolumn{
  % Main space
  % \textbf{Reproducible GPU environments} with Pixi, Linux container images, and CHTC resources make complex scientific workflows \textbf{easy to run}
  \textbf{Reproducible GPU workflows} with Pixi, Linux container images, and CHTC resources make complex hardware-accelerated scientific tasks \textbf{easy to run at scale}
 }{
  % Bottom space

\begin{center}
  % QR code
  \qrcode{qr_code.pdf}{smartphoneWhite}{
   \textbf{Take a picture} to get more
   \\information, examples, and this poster\\
   \href{https://github.com/matthewfeickert/poster-chtc-researcher-forum-2026}{github.com/matthewfeickert/poster-chtc-researcher-forum-2026}
  }
\end{center}
}

}{
 % LEFT COLUMN
 \title{Accelerating hardware-accelerated science}
 \vspace{0.5em}
 \author{{\href{https://www.matthewfeickert.com/}{Matthew Feickert}$^{1}$}}
 \vspace{1.0em}
 \institution{1 University of Wisconsin--Madison Data Science Institute}
 %
 \section{Hardware-accelerated science}
 Machine learning is becoming a standard part of the modern scientist's toolkit.
 As a result, the use of \textbf{hardware accelerators} (e.g. GPUs) in scientific software workflows is increasingly common.
 Historically, GPU-enabled software could be difficult to work with and required expertise to deploy.

 \vspace{1.0em}
 \textbf{Reproducible environments with Pixi}
 \vspace{1.0em}
  \begin{center}
   \href{https://pixi.prefix.dev/}{\includegraphics[width=0.5\textwidth]{tools_pixi.pdf}}
  \end{center}
 \begin{itemize}
  \item Declaratively create \textbf{fully reproducible locked environments} with \textbf{conda} packages (ANY language) and \textbf{Python} packages
  \item Locking: automatic and Rust-fast
  \item Define host-specific expectations of hardware resources (e.g. CUDA version)
  \item Resolve environments for \textbf{multiple different platforms} than solve machine (e.g. solve on Mac laptop, run on Linux machine with CUDA)
 \end{itemize}

 \vspace{1.0em}
 \textbf{Linux container distributions}

 \vspace{1.0em}
 \begin{minipage}{0.50\textwidth}
  \begin{center}
   \href{https://www.docker.com/}{\includegraphics[width=0.8\textwidth]{docker-logo-blue.pdf}}
  \end{center}
 \end{minipage}%
 \hfill
 \begin{minipage}{0.50\textwidth}
  \begin{center}
   \href{https://apptainer.org/}{\includegraphics[width=0.8\textwidth]{apptainer_logo.pdf}}
  \end{center}
 \end{minipage}%
 \vspace{1.0em}

 \begin{itemize}
  \item Deploy software environments to worker nodes through Docker and Apptainer \textbf{container images}
  \item With Pixi environment, Dockerfile or Apptainer definition files become a \textbf{simple template}
  \item Changes to environment can trigger rebuilds and redeployments of container images
  \item Fast launching with CHTC's \texttt{/staging} and OSDF
 \end{itemize}
}{
 % RIGHT COLUMN

 \textbf{HTCondor containerized GPU jobs}

 \vspace{1.0em}
 \begin{minipage}{0.50\textwidth}
  \begin{center}
   \href{https://chtc.cs.wisc.edu/}{\includegraphics[width=0.85\textwidth]{chtc_logo.pdf}}
  \end{center}
 \end{minipage}%
 \hfill
 \begin{minipage}{0.50\textwidth}
  \begin{center}
   \href{https://htcondor.org/}{\includegraphics[width=\textwidth]{htcondor_logo.pdf}}
  \end{center}
 \end{minipage}%
 \vspace{1.0em}
 \begin{itemize}
  \item CHTC's HTCondor deployment support \textbf{container universe}
  \item Specific control over \textbf{GPU and CUDA capabilities} through submit description files
  \item Parallel tasks can scaleout across CHTC GPU nodes
 \end{itemize}

 \subsection{Implementation}

 \vspace{2.0em}
 \begin{center}
  \includegraphics[width=\textwidth]{implimentation.png}
 \end{center}
 \vspace{2.0em}

 \subsection{Applications in DSI driven science}
 Data scientists and researchers at the \textbf{University of Wisconsin--Madison Data Science Institute} are using this workflow pattern across:
 \begin{itemize}
  \item High Energy Physics applications of simulation-based inference
  \item High-Resolution Rapid Refresh (HRRR) weather forecasting
  \item Large-scale simulation of protein motions using molecular dynamics
 \end{itemize}

 \subsection{More information}
 \begin{itemize}
  \item Matthew Feickert, Ruben Arts, John Kirkham, \emph{Reproducible Machine Learning Workflows for
Scientists with Pixi}, \\\textbf{DOI: \href{https://doi.org/10.25080/nwuf8465}{10.25080/nwuf8465}}
  \item \href{https://github.com/CHTC/templates-GPUs}{\textbf{github.com/CHTC/templates-GPUs}}
  \item \href{https://github.com/CHTC/recipes}{\textbf{github.com/CHTC/recipes}}
 \end{itemize}

\vfill
\begin{center}
  \href{https://dsi.wisc.edu/}{\includegraphics[width=\textwidth]{DSI_logo.png}}
\end{center}
}
\end{document}
